\documentclass[a4paper]{article}
\usepackage[pdftex]{graphicx}
\usepackage[utf8]{inputenc}
\usepackage{enumerate}
\usepackage{amssymb}
\usepackage{icomma}
\usepackage{colortbl}
\usepackage{siunitx}
\sisetup{locale=DE} 
\usepackage{tikz}
\usepackage{href-ul}
\hypersetup{
	colorlinks=true,
	linkcolor=blue,
	urlcolor=blue}
\usepackage{geometry}
\geometry{a4paper, top=15mm, left=15mm, right=15mm, bottom=15mm,
	headsep=10mm, footskip=12mm}

\begin{document}
\begin{tikzpicture}(10,3)
	\draw[ultra thick](2,0) --(10,0) -- (10,1.3) --(2,1.3) -- (2,0);
	\draw[fill=black](2,0)-- (0,.6) -- (2,.6) -- (2,0);
	\node at (6,.6) {\large Lösungsblatt von \href{https://www.okuyakl.de}{www.okuyakl.de}};
\end{tikzpicture}
\vspace{0.5 cm}

\begin{enumerate}[1.]
{\bf \item Berechne die fehlenden Größen}


	\renewcommand{\arraystretch}{2}
	\begin{tabular}{|>{\columncolor[gray]{.8}}p{1.5 cm}|p{1.5 cm}|p{1.5 cm}|p{1.5 cm}|p{1.5 cm}|p{1.5 cm}|}
		\hline
		\rowcolor[gray]{.8}	&a) &b) &c) & d) & e) \\
		\hline
		$a$ & $\SI{4,2}{\centi\meter}$ &$\SI{2,8}{\meter}$ &$\SI{36}{\centi\meter}$ &$\SI{1,2}{\kilo\meter}$ & $a$\\
		\hline
		$b$ &$\SI{12,2}{\centi\meter}$ & $\SI{27}{\meter}$&$\SI{3,5}{\deci\meter}$ &$\SI{212}{\meter}$ &$\sqrt{3}\cdot a$ \\
		\hline
		$\alpha$ &$19^\circ$ &$6^\circ$ &$46^\circ$ & $10^\circ$ & $30^\circ$\\
		\hline
		$\beta$ &$71^\circ$ &$84^\circ$ &$44^\circ$ &$8 \cdot \alpha$ &$60^\circ$ \\
		\hline
	\end{tabular}

{\bf \item Aufgabe}
	\begin{enumerate}[a)]
		\item $$\alpha=32^\circ \qquad \beta=58^\circ$$
		\item $$\gamma=2\cdot32^\circ = 64^\circ$$
	\end{enumerate} 


{\bf \item Aufgabe}

Wir messen: $$e=\SI{2,3}{\centi\meter} \qquad f=\SI{4,9}{\centi\meter}$$

$$\tan {\alpha\over 2} = {0,5 \cdot e \over 0,5 \cdot f} \Rightarrow {\alpha = 50^\circ}$$\\
$$\beta = 180^\circ -\alpha = 130^\circ$$

 
{\bf \item Aufgabe}
\begin{enumerate}[a)]
	\item $$r={\SI{3500}{\kilo\meter}\over \tan{0,5^\circ}}\approx \SI{400000}{\kilo\meter}$$
	\item $$d_S=\SI{150e6}{\kilo\meter}\cdot \tan{0,5}=\SI{1,3e6}{\kilo\meter}$$
	\item $$1~\rm{pc} = {\SI{300e6}{\kilo\meter} \over \tan{{1\over 3600}^\circ}}=\SI{3,1e13}{\kilo\meter}$$
	\item $$1~\rm{pc} ={\SI{3,1e13}{\kilo\meter} \over \SI{9,5e12}{\kilo\meter}}=3,3 ~\rm{Lichtjahre}$$
\end{enumerate}  

{\bf \item Aufgabe}
\begin{enumerate}[a)]
	\item $$d= {\SI{1731}{\meter} \over \tan{8^\circ}}=\SI{12}{\kilo\meter}$$
	\item $$ \tan \alpha = 0,05 \Rightarrow \alpha = 2,9^\circ$$
\end{enumerate} 



{\bf \item Aufgabe}

\begin{enumerate}[a)]
	\item $$\tan \alpha={\SI{28}{\centi\meter}\over\SI{70}{\centi\meter}} \Rightarrow \alpha= 21,8^\circ$$ 
	\item $$h=\SI{30}{\meter} \cdot \tan{21,8^\circ}=\SI{12}{\meter}$$
	\item $$h={\SI{24}{\centi\meter}\over\SI{70}{\centi\meter}}\cdot \SI{30}{\meter}= \SI{12}{\meter}$$ 
\end{enumerate} 
 

	{\bf \item Aufgabe}
\begin{enumerate}[a)]
	\item 
	\begin{minipage}{0.5\textwidth}
			\includegraphics[width=7 cm]{gera100}	
	\end{minipage}
	\hfill
	\begin{minipage}{0.5\textwidth}
		\item $$\alpha_g=27^\circ \qquad \alpha_h=-72^\circ \qquad \alpha_f=-22^\circ$$
		\item $$\angle (g,h) = 81^\circ \qquad \angle (g,f) = 49^\circ \qquad \angle (h,f)=50^\circ$$  
	\end{minipage}
\end{enumerate} 

{\bf \item Aufgabe}

\begin{enumerate}[a)]
\item
$$A_1(2|2)\qquad B_1(4|2)$$

\item Wir berechnen die Streckenlängen:
$$\overline{A_1B_1}=\left|{4-2 \choose 2-2}\right|=\left|{2 \choose 0}\right|= 2$$
$$\overline{B_1C_1}=\left|{4-4 \choose 4-2}\right|=\left|{0 \choose 2}\right|= 2$$
Damit haben wir ein gleichschenklig-rechtwinkliges Dreieck und es gilt: $$\tan\alpha=1 \quad \Rightarrow \quad \alpha=45^\circ$$

\item Allgemein ist:
$$\overline{A_nB_n}=\left|{4-x \choose 0}\right|= 4 - x $$
und
$$\overline{B_nC_n}=\left|{0 \choose 4+0,5x-3}\right|= 0,5x + 1$$
$$\tan\alpha={\overline{B_nC_n} \over \overline{A_nB_n} } = {0,5x+1 \over 4-x}$$

\item 
$$
\renewcommand{\arraystretch}{2}
\begin{array}{rcl}
\tan 14,04^\circ = 0,25 &=& {0,5x+1 \over 4-x}\\
0,25\cdot (4-x) &=& 0,5x+1\\
1 -0,25x &=& 0,5x +1 \\
0 &=& 0,75x \\
x&=& 0\\
A_2(0|3) & & B_2(4|3)
\end{array}
$$

\end{enumerate} 


\begin{center}
	\includegraphics[width=7 cm]{../../viecher/endcomic.pdf}
	
	Hier geht es zurück zum \href{https://www.okuyakl.de/math/m10tanL100/aa100.pdf}{Aufgabenblatt}
\end{center}

\end{enumerate}
\end{document}

